\subsection{Eliminasi Gauss}

Pekerjaan yang kita lakukan pada bagian sebelumnya akan selalu menemukan solusi sistem. Pada bagian ini, kita akan mengkaji cara yang lebih sederhana untuk
menemukan solusi. Pertama, kita akan merepresentasikan sistem linear dengan \textbf{matriks diperbesar.} \index{matriks!matriks diperbesar}
Suatu \textbf{matriks} \index{matriks} hanyalah susunan bilangan berbentuk persegi panjang. Ukuran atau dimensi suatu matriks \index{matriks!dimensi} didefinisikan sebagai $m\times n$, dengan $m$ menyatakan
banyaknya baris dan $n$ menyatakan banyaknya kolom. Untuk menyusun matriks diperbesar
dari suatu sistem linear, kita membentuk \textbf{matriks koefisien} \index{matriks!matriks koefisien}dari koefisien variabel-variabel dalam sistem, serta
\textbf{matriks konstanta} dari konstanta-konstantanya. Koefisien dari satu persamaan dalam sistem membentuk satu baris matriks diperbesar.

Sebagai contoh, perhatikan sistem linear pada Contoh \ref{exa:solvingasystemwithelementaryops}
\begin{equation*}
\begin{array}{ccccccc}
x&+&3y&+&6z&=&25 \\
2x&+&7y&+&14z&=&58 \\
&&2y&+&5z&=&19
\end{array}
\end{equation*}
Sistem ini dapat ditulis sebagai matriks diperbesar, sebagai berikut \:
\begin{equation*}
\leftB 
\begin{array}{rrr|r}
1 & 3 & 6 & 25 \\
2 & 7 & 14 & 58 \\
0 & 2 & 5 & 19
\end{array}
\rightB 
\end{equation*}

Perhatikan bahwa matriks tersebut memuat informasi yang persis sama dengan sistem semula. Di sini dipahami
bahwa kolom pertama memuat koefisien $x$ dalam setiap persamaan secara berurutan,
$\leftB
\begin{array}{r}
1 \\
2 \\
0
\end{array}
\rightB .$ Dengan cara serupa, kita membentuk kolom dari koefisien $y$ dalam setiap persamaan, $\leftB
\begin{array}{r}
3 \\
7 \\
2
\end{array}
\rightB $, serta kolom dari koefisien $z$ dalam setiap persamaan, $\leftB
\begin{array}{r}
6 \\
14 \\
5
\end{array}
\rightB .$ Untuk sistem dengan lebih dari tiga variabel, kita melanjutkan cara ini dengan membentuk satu kolom untuk setiap variabel.
Demikian pula, untuk sistem dengan kurang dari tiga variabel, kita cukup membentuk satu kolom untuk setiap variabel.

Terakhir, kita membentuk kolom dari konstanta-konstanta persamaan,
$\leftB
\begin{array}{r}
25\\
58\\
19
\end{array}
\rightB .$ 

Baris-baris matriks diperbesar bersesuaian dengan persamaan-persamaan dalam sistem. Sebagai contoh, baris
teratas dalam matriks diperbesar,
$\leftB
\begin{array}{rrrrr}
1 & 3 & 6 & | & 25
\end{array}
\rightB$
 bersesuaian dengan persamaan
\begin{equation*}
x+3y+6z=25.
\end{equation*}

Perhatikan definisi berikut.

\begin{definition}{Matriks Diperbesar Suatu Sistem Linear}\label{def:augmentedmatrix}
Untuk sistem linear berbentuk
\begin{equation*}
\begin{array}{c}
a_{11}x_{1}+\cdots +a_{1n}x_{n}=b_{1} \\
\vdots \\
a_{m1}x_{1}+\cdots +a_{mn}x_{n}=b_{m}
\end{array}
\end{equation*}
dengan $x_{i}$ berupa variabel serta $a_{ij}$ dan $b_{i}$ berupa konstanta,
matriks diperbesar \index{matriks!matriks diperbesar} dari sistem ini diberikan oleh
\begin{equation*}
\leftB
\begin{array}{rrr|r}
a_{11} & \cdots & a_{1n} &  b_{1} \\
\vdots &  & \vdots &  \vdots \\
a_{m1} & \cdots & a_{mn} &  b_{m}
\end{array}
\rightB 
\end{equation*}
\end{definition}

Sekarang, perhatikan operasi elementer dalam konteks matriks diperbesar. Operasi elementer pada Definisi \ref{def:elementaryoperations}
dapat digunakan pada baris sebagaimana sebelumnya kita menggunakannya pada persamaan. Perubahan pada suatu sistem persamaan akibat
operasi elementer ekuivalen dengan perubahan pada matriks diperbesar akibat operasi baris yang bersesuaian. Perhatikan bahwa Teorema \ref{thm:elementaryoperationsandsolns} menyiratkan bahwa
setiap operasi baris elementer yang digunakan pada matriks diperbesar tidak akan mengubah solusi sistem persamaan yang bersesuaian.
Sekarang kita mendefinisikan operasi baris elementer secara formal. Operasi ini merupakan {\em alat utama \em}yang akan kita gunakan untuk menemukan solusi sistem persamaan.

\begin{definition}{Operasi Baris Elementer}\label{def:rowoperations}
\textbf{Operasi baris elementer} (juga dikenal sebagai \textbf{operasi baris}) terdiri atas operasi-operasi berikut \index{operasi baris} \index{operasi baris elementer}

\begin{enumerate}
\item Tukarkan dua baris.

\item Kalikan suatu baris dengan bilangan tak nol.

\item Ganti suatu baris dengan baris tersebut ditambah sembarang kelipatan dari baris lain.
\end{enumerate}
\end{definition}

Ingat kembali cara kita menyelesaikan Contoh \ref{exa:solvingasystemwithelementaryops}.
Kita dapat melakukan langkah-langkah yang persis sama seperti di atas, tetapi sekarang dalam konteks matriks diperbesar dan dengan menggunakan operasi baris.
Matriks diperbesar sistem ini adalah
\begin{equation*}
\leftB 
\begin{array}{rrr|r}
1 & 3 & 6 &  25 \\
2 & 7 & 14 &  58 \\
0 & 2 & 5 &  19
\end{array}
\rightB 
\end{equation*}
Jadi, langkah pertama dalam
menyelesaikan sistem yang diberikan oleh \ref{solvingasystem1} adalah mengambil $\left( -2\right) $ kali baris pertama
matriks diperbesar dan menambahkannya ke baris kedua,
\begin{equation*}
\leftB
\begin{array}{rrr|r}
1 & 3 & 6 & 25 \\
0 & 1 & 2 & 8 \\
0 & 2 & 5 & 19
\end{array}
\rightB 
\end{equation*}
Perhatikan bahwa hal ini bersesuaian dengan \ref{solvingasystem2}. Selanjutnya, ambil $\left( -2\right) $
kali baris kedua dan tambahkan ke baris ketiga,
\begin{equation*}
\leftB 
\begin{array}{rrr|r}
1 & 3 & 6 & 25 \\
0 & 1 & 2 & 8 \\
0 & 0 & 1 & 3
\end{array}
\rightB
\end{equation*}
Matriks diperbesar ini bersesuaian dengan sistem
\begin{equation*}
\begin{array}{ccccccc}
x&+&3y&+&6z&=&25 \\
&&y&+&2z&=&8 \\
&&&&z&=&3
\end{array}
\end{equation*}
yang sama dengan \ref{solvingasystem3}. Dengan substitusi balik, Anda memperoleh solusi
$x=1,y=2,$ dan $z=3.$

Melalui prosedur operasi baris yang
sistematis, kita dapat menyederhanakan matriks diperbesar dan membawanya ke \textbf{\ef} atau \textbf{\rref}, yang akan didefinisikan berikutnya. Bentuk-bentuk
ini digunakan untuk menemukan solusi sistem persamaan yang bersesuaian dengan matriks diperbesar.

Dalam definisi-definisi berikut, istilah
\textbf{entri utama}
\index{entri utama}merujuk pada
 entri tak nol pertama dalam suatu baris ketika baris tersebut dibaca
dari kiri ke kanan.

\begin{definition}{\EF}\label{def:echelonform}
Suatu matriks diperbesar berada
\index{\ef} dalam \textbf{\ef} jika

\begin{enumerate}
\item Semua baris tak nol berada di atas setiap baris nol.

\item Setiap entri utama suatu baris berada dalam kolom di sebelah kanan entri utama setiap baris di atasnya.

%%Definition using original Kuttler definition ends here

%%Definition using Style B definition of \ref
\item Setiap entri utama suatu baris sama dengan $1$.

\end{enumerate}
\end{definition}

%%For Style B:
Kita juga membahas bentuk tereduksi lain dari matriks diperbesar yang memiliki satu syarat tambahan.

%%For Original Kuttler:
%% We also consider another reduced form of the augmented matrix which has further conditions.

\begin{definition}{\RREF}\label{def:rref}
Suatu matriks
\index{\rref} diperbesar berada dalam \textbf{\rref } jika

\begin{enumerate}
\item Semua baris tak nol berada di atas setiap baris nol.

\item Setiap entri utama suatu baris berada dalam kolom di sebelah kanan entri utama
setiap baris di atasnya.

\item Setiap entri utama suatu baris sama dengan $1$.

\item Semua entri dalam suatu kolom yang berada di atas dan di bawah entri utama bernilai nol.

\end{enumerate}
\end{definition}

%%For Style B definition:
Perhatikan bahwa tiga syarat pertama pada matriks \rref \;sama dengan syarat untuk \ef.

%%For original Kuttler definition:
%%Notice that the first two conditions on a \rref  \; matrix are the same as those for \ef.

Jadi, setiap matriks \rref \;juga berada dalam \ef. Kebalikannya belum tentu benar;
 kita tidak dapat menganggap bahwa setiap matriks dalam \ef \;juga berada dalam \rref. Namun, sering kali
\ef \;sudah cukup untuk memberikan informasi mengenai solusi suatu sistem.

Contoh-contoh berikut memperlihatkan matriks dalam berbagai bentuk tersebut. Sebagai latihan, luangkan waktu
untuk memverifikasi dengan cermat bahwa matriks-matriks tersebut berada dalam bentuk yang dinyatakan.

\begin{example}{Tidak Berada dalam \EF}\label{exa:matricesnotechelon}
Matriks-matriks diperbesar berikut tidak berada dalam \ef \;(dan karena itu juga tidak berada dalam \rref).

\begin{equation*}
\leftB
\begin{array}{rrr|r}
0 & 0 & 0 & 0 \\
1 & 2 & 3 & 3 \\
0 & 1 & 0 & 2 \\
0 & 0 & 0 & 1 \\
0 & 0 & 0 & 0
\end{array}
\rightB ,\leftB
\begin{array}{rr|r}
1 & 2 & 3 \\
2 & 4 & -6 \\
4 & 0 & 7
\end{array}
\rightB ,\leftB
\begin{array}{rrr|r}
0 & 2 & 3 & 3 \\
1 & 5 & 0 & 2 \\
7 & 5 & 0 & 1 \\
0 & 0 & 1 & 0
\end{array}
\rightB 
\end{equation*}
\end{example}

\begin{example}{Matriks dalam \EF}\label{exa:matricesnotrref}

%%For Original Kuttler definition, can change these matrices to not have leading 1s.

Matriks-matriks diperbesar berikut berada dalam \ef, tetapi tidak berada dalam \rref.
\begin{equation*}
\leftB
\begin{array}{rrrrr|r}
1 & 0 & 6 & 5 & 8 & 2 \\
0 & 0 & 1 & 2 & 7 & 3 \\
0 & 0 & 0 & 0 & 1 & 1 \\
0 & 0 & 0 & 0 & 0 & 0
\end{array}
\rightB ,\leftB
\begin{array}{rrr|r}
1 & 3 & 5 & 4 \\
0 & 1 & 0 & 7 \\
0 & 0 & 1 & 0 \\
0 & 0 & 0 & 1 \\
0 & 0 & 0 & 0
\end{array}
\rightB, \leftB
\begin{array}{rrr|r}
1 & 0 & 6 & 0 \\
0 & 1 & 4 & 0 \\
0 & 0 & 1 & 0 \\
0 & 0 & 0 & 0
\end{array}
\rightB
\end{equation*}
\end{example}

Perhatikan bahwa kita dapat menerapkan operasi baris lebih lanjut pada matriks-matriks ini untuk membawanya ke \rref. Luangkan waktu untuk mencobanya sendiri.
Perhatikan matriks-matriks berikut, yang berada dalam \rref.

\begin{example}{Matriks dalam \RREF}\label{exa:rrefmatrices}
Matriks-matriks diperbesar berikut berada dalam \rref.
\begin{equation*}
\leftB
\begin{array}{rrrrr|r}
1 & 0 & 0 & 5 & 0 & 0 \\
0 & 0 & 1 & 2 & 0 & 0 \\
0 & 0 & 0 & 0 & 1 & 1 \\
0 & 0 & 0 & 0 & 0 & 0
\end{array}
\rightB ,\leftB
\begin{array}{rrr|r}
1 & 0 & 0 & 0 \\
0 & 1 & 0 & 0 \\
0 & 0 & 1 & 0 \\
0 & 0 & 0 & 1 \\
0 & 0 & 0 & 0
\end{array}
\rightB , \leftB
\begin{array}{rrr|r}
1 & 0 & 0 & 4 \\
0 & 1 & 0 & 3 \\
0 & 0 & 1 & 2 
\end{array}
\rightB
\end{equation*}
\end{example}

Salah satu kegunaan \ef \;suatu matriks adalah untuk menentukan posisi pivot dan kolom pivot matriks tersebut.

\begin{definition}{Posisi Pivot dan Kolom Pivot}\label{def:pivotpositionandcolumn}
Suatu \textbf{posisi pivot}
\index{posisi pivot}dalam matriks adalah letak entri utama dalam \ef{} suatu matriks.
Suatu \textbf{kolom pivot}
\index{kolom pivot} adalah kolom yang memuat posisi pivot.
\end{definition}

Sebagai contoh, perhatikan yang berikut.

\begin{example}{Posisi Pivot}\label{exa:pivotposition}
Misalkan
\begin{equation*}
A=\leftB
\begin{array}{rrr|r}
1 & 2 & 3 & 4 \\
3 & 2 & 1 & 6 \\
4 & 4 & 4 & 10
\end{array}
\rightB
\end{equation*}
Di manakah posisi pivot
dan kolom pivot matriks diperbesar $A$?
\end{example}

\begin{solution} 
%%Solution as per Style B definition of row-echelon form
\ef \;dari matriks ini adalah
\begin{equation*}
\leftB
\begin{array}{rrr|r}
1 & 2 & 3 & 4 \\
0 & 1 & 2 & \vspace{0.05in}\frac{3}{2} \\
0 & 0 & 0 & 0
\end{array}
\rightB
\end{equation*}

%%Solution as per Kuttler's original definition of echelon form
%% The \ef \;of this matrix is
%\begin{equation*}
%\leftB
%\begin{array}{rrr|r}
%%1 & 2 & 3 & 4 \\
%%0 & -4& -8 & -6 \\
%%0 & 0 & 0 & 0
%\end{array}
%\rightB
%\end{equation*} 

Inilah yang kita perlukan dalam contoh ini, tetapi perhatikan bahwa matriks tersebut tidak berada dalam \rref.

Untuk menentukan posisi pivot dalam matriks semula, kita mencari entri utama dalam \ef \hspace{0.5mm} matriks tersebut.
Di sini, entri pada baris pertama dan kolom pertama serta entri pada baris kedua dan kolom kedua
merupakan entri utama. Jadi, letak-letak tersebut merupakan posisi pivot. Kita menandai posisi pivot pada matriks semula,
sebagai berikut:\
\begin{equation*}
\leftB
\begin{array}{rrr|r}
\fbox{1} & 2 & 3 & 4 \\
3 & \fbox{2} & 1 & 6 \\
4 & 4 & 4 & 10
\end{array}
\rightB
\end{equation*}
Jadi, kolom pivot dalam matriks tersebut adalah dua kolom pertama.
\end{solution}

Berikut adalah algoritme untuk membawa suatu matriks
ke \ef \;dan \rref. Sebagai latihan, Anda dapat menggunakan algoritme ini untuk membawa sendiri matriks di atas ke \ef \;atau \rref.

\begin{algorithm}
\caption{Algoritme \RREF}
\label{algo:rrefalgorithm}

Algoritme ini memberikan metode untuk menggunakan operasi baris guna membawa suatu matriks ke
\rref-nya \index{\rref!algoritme}.
Kita memulai dengan matriks dalam bentuk semula.

\begin{enumerate}
\item Mulai dari kiri, temukan kolom tak nol pertama. Kolom ini merupakan kolom pivot
pertama, dan posisi pada bagian atas kolom ini merupakan posisi pivot
pertama. Tukarkan baris jika perlu untuk menempatkan bilangan tak nol pada
posisi pivot pertama.

\item Gunakan operasi baris untuk membuat entri-entri di bawah posisi pivot
pertama (dalam kolom pivot pertama) sama dengan nol.

%%Algorithm for the Style B definition of row-echelon form
\item Dengan mengabaikan baris yang memuat posisi pivot pertama, ulangi langkah 1 dan 2 pada baris-baris yang tersisa.
Ulangi proses tersebut sampai tidak ada lagi baris
yang perlu diubah.

\item Bagi setiap baris tak nol dengan nilai entri utamanya, sehingga entri utama menjadi $1$. Matriks tersebut kemudian berada dalam \ef.

\end{enumerate}
Langkah berikut akan membawa matriks dari \ef \;ke \rref.

\begin{enumerate}
\setcounter{enumi}{4}

\item Dengan bergerak dari kanan ke kiri, gunakan operasi baris untuk menghasilkan nol pada entri-entri kolom pivot yang berada di atas posisi pivot. Hasilnya
berupa matriks dalam \rref.
\end{enumerate}

%%Algorithm for the original Kuttler definition of echelon form
%\item Ignoring the row containing the first pivot position, repeat steps 1 and 2 with the remaining rows. Repeat the process until there 
%% are not more rows to modify. The matrix will then be in \ef.
%\suspend{enumerate}
%%The following steps will carry the matrix from \ef \;to \rref.
%\resume{enumerate}
%\item Moving from right to left, use row operations to create zeros in the entries of the pivot columns which are above the pivot positions.
%\item Divide each nonzero row by the value of the leading entry, so that the leading entry becomes $1$. The result will be 
%% a matrix in \rref.
%\end{enumerate}
\end{algorithm}

Paling sering, kita akan menerapkan algoritme ini pada matriks diperbesar untuk menemukan solusi
suatu sistem persamaan linear. Namun, kita dapat menggunakan algoritme ini untuk menghitung \rref \;dari sembarang matriks, yang dapat berguna dalam penerapan lain.

Perhatikan contoh Algoritme \ref{algo:rrefalgorithm} berikut.

\begin{example}{Menentukan \EF{} dan \\ \RREF{} Suatu Matriks}\label{exa:findingechelonandrref}
Misalkan
\begin{equation*}
A = \leftB
\begin{array}{rrr}
0 & -5 & -4 \\
1 & 4 & 3 \\
5 & 10 & 7
\end{array}
\rightB
\end{equation*}
Temukan \ef \;dari $A$. Kemudian selesaikan
prosesnya hingga $A$ berada dalam \rref.
\end{example}

\begin{solution} Dalam mengerjakan contoh ini, kita akan menggunakan langkah-langkah yang diuraikan dalam Algoritme \ref{algo:rrefalgorithm}.
\begin{enumerate}

\item Kolom pivot pertama adalah kolom pertama matriks karena kolom ini merupakan kolom tak nol pertama dari kiri.
Jadi, posisi pivot pertama berada pada baris pertama dan kolom pertama. Tukarkan kedua baris pertama untuk
memperoleh entri tak nol pada posisi pivot pertama, yang ditandai dengan kotak di bawah.
\begin{equation*}
\leftB
\begin{array}{rrr}
\fbox{1} & 4 & 3 \\
0 & -5 & -4 \\
5 & 10 & 7 
\end{array}
\rightB
\end{equation*}

\item Langkah kedua menghasilkan nol pada entri-entri di bawah posisi pivot pertama. Entri pertama pada baris kedua
sudah bernilai nol. Kita hanya perlu mengurangkan $5$ kali baris pertama dari baris ketiga.
Matriks yang dihasilkan adalah
\begin{equation*}
\leftB
\begin{array}{rrr}
1 & 4 & 3 \\
0 & -5 & -4 \\
0 & 10 & 8 
\end{array}
\rightB
\end{equation*} 

%%Solution using Style B definition of \ef

\item Sekarang abaikan baris teratas. Terapkan langkah $1$ dan $2$ pada matriks yang lebih kecil
\begin{equation*}
\leftB
\begin{array}{rr}
-5 & -4\\
10 & 8 
\end{array}
\rightB
\end{equation*}
Dalam matriks ini, kolom pertama merupakan kolom pivot dan $-5$ berada pada posisi pivot pertama. Oleh karena itu, kita perlu menghasilkan nol di bawahnya.
Untuk melakukannya, tambahkan $2$ kali baris pertama (dari matriks ini) ke baris kedua. Matriks yang dihasilkan adalah
\begin{equation*}
\leftB
\begin{array}{rr}
-5 & -4\\
0 & 0
\end{array}
\rightB
\end{equation*}
Matriks semula kini tampak seperti
\begin{equation*}
\leftB
\begin{array}{rrr}
1 & 4 & 3 \\
0 & -5 & -4 \\
0 & 0 & 0 
\end{array}
\rightB
\end{equation*}
Kita dapat melihat bahwa tidak ada lagi baris yang perlu diubah.

\item Sekarang kita perlu menghasilkan $1$ utama pada setiap baris. Baris pertama sudah memiliki $1$ utama, sehingga tidak diperlukan tindakan di sini.
Bagi baris kedua dengan $-5$ untuk menghasilkan $1$ utama. Matriks yang dihasilkan adalah
\begin{equation*}
\leftB
\begin{array}{rrr}
1 & 4 & 3 \\
0 & 1 & \vspace{0.05in}\frac{4}{5} \\
0 & 0 & 0
\end{array}
\rightB
\end{equation*}
Matriks ini kini berada dalam \ef.

\item Sekarang hasilkan nol pada entri-entri di atas posisi pivot dalam setiap kolom untuk membawa matriks ini sepenuhnya ke \rref.
Perhatikan bahwa tidak terdapat posisi pivot pada kolom ketiga, sehingga kita tidak perlu menghasilkan nol dalam kolom tersebut. Kolom tempat kita
perlu menghasilkan nol adalah kolom kedua. Untuk melakukannya, kurangkan $4$ kali baris kedua dari baris pertama. Matriks yang dihasilkan adalah
\begin{equation*}
\leftB
\begin{array}{rrr}
1 & 0 & -\vspace{0.05in}\frac{1}{5} \\
0 & 1 & \vspace{0.05in}\frac{4}{5} \\
0 & 0 & 0
\end{array}
\rightB
\end{equation*}
\end{enumerate}
Matriks ini kini berada dalam \rref.
\end{solution}

%%Solution using the original Kuttler definition of \ef
%\item Now ignore the top row.  Apply steps $1$ and $2$ to the smaller matrix
%\begin{equation*}
%\leftB
%\begin{array}{rr}
%%-5 & -4\\
%%10 & 8 
%\end{array}
%\rightB.
%\end{equation*}
%%In this matrix, $-5$ is in the pivot position. Therefore, we need to create a zero below it.
%%To do this, add $2$ times the first row (of this matrix) to the second. The resulting matrix is
%\begin{equation*}
%\leftB
%\begin{array}{rr}
%%-5 & 4\\
%%0 & 0
%\end{array}
%\rightB.
%\end{equation*}
%%Our original matrix now looks like
%\begin{equation*}
%\leftB
%\begin{array}{rrr}
%%1 & 4 & 3 \\
%%0 & -5 & -4 \\
%%0 & 0 & 0 
%\end{array}
%\rightB
%\end{equation*}
%%We can see that there are no more rows to modify. Therefore, this matrix is in \ef.

%\item Now, we wish to continue carrying this matrix all the way to \rref.
%%First, we want to create zeros in the entries above pivot positions in each column. 
%%Notice that there is no pivot position in the third column so we do not need to create any zeros in this column! The column which we 
%%need to create zeros in is the second. To do so, add $\frac{4}{5}$ times the second row from the first row. The resulting matrix is
%\begin{equation*}
%\leftB
%\begin{array}{rrr}
%%1 & 0 & -\vspace{0.05in}\frac{1}{5} \\
%%0 & -5 & -4 \\
%%0 & 0 & 0
%\end{array}
%\rightB
%\end{equation*}

%\item Finally, we need to create leading $1$s in each row. The first row already has a leading $1$ so no work is needed here.
%%Divide the second row by $-5$ to create a leading $1$. The resulting matrix is
%\begin{equation*}
%\leftB
%\begin{array}{rrr}
%%1 & 0 & -\vspace{0.05in}\frac{1}{5} \\
%%0 & 1 & \vspace{0.05in}\frac{4}{5} \\
%%0 & 0 & 0
%\end{array}
%\rightB
%\end{equation*}
%\end{enumerate}
%%This matrix is now in \rref.
%\end{solution}

Algoritme di atas memberi Anda cara sederhana untuk memperoleh \ef \;dan \rref \;suatu matriks.
Gagasan utamanya adalah melakukan operasi baris sedemikian rupa sehingga diperoleh
matriks dalam \ef \;atau \rref. Proses ini penting karena matriks yang dihasilkan
memungkinkan Anda mendeskripsikan
solusi sistem persamaan linear yang bersesuaian secara bermakna.

Pada contoh berikutnya, kita akan melihat cara menyelesaikan sistem persamaan menggunakan matriks diperbesar yang bersesuaian.

\begin{example}{Menemukan Solusi Suatu Sistem}\label{exa:systemwithnosoln}
Berikan solusi lengkap sistem persamaan berikut \:
\begin{equation*}
\begin{array}{ccccccccc}
2x&+&4y&-&3z&=&-1\\
5x&+&10y&-&7z&=&-2\\
3x&+&6y&+&5z&=&9
\end{array}
\end{equation*}
\end{example}

\begin{solution} Matriks diperbesar sistem ini adalah
\begin{equation*}
\leftB
\begin{array}{rrr|r}
2 & 4  & -3 & -1 \\
5 & 10 & -7 & -2 \\
3 & 6  &  5 &  9
\end{array}
\rightB
\end{equation*}

Untuk menemukan solusi sistem ini, kita ingin membawa matriks diperbesar ke \rref.
Kita akan melakukannya dengan menggunakan Algoritme \ref{algo:rrefalgorithm}. Perhatikan bahwa
kolom pertama tidak nol, sehingga kolom ini merupakan kolom pivot pertama. Entri pertama pada baris pertama, yaitu $2$, merupakan
entri utama pertama dan berada pada posisi pivot pertama. Kita akan menggunakan operasi baris
untuk menghasilkan nol pada entri-entri di bawah $2$. Pertama, ganti baris kedua dengan $-5$ kali baris pertama ditambah $2$ kali baris kedua. Hal ini menghasilkan
\begin{equation*}
\leftB
\begin{array}{rrr|r}
2 & 4 & -3 & -1 \\
0 & 0 & 1 & 1 \\
3 & 6 & 5 & 9
\end{array}
\rightB
\end{equation*}
Sekarang, ganti baris ketiga dengan $-3$ kali baris pertama
ditambah $2$ kali baris ketiga. Hal ini menghasilkan
\begin{equation*}
\leftB
\begin{array}{rrr|r}
2 & 4 & -3 &  -1 \\
0 & 0 & 1 & 1 \\
0 & 0 & 1 & 21
\end{array}
\rightB 
\end{equation*}
Sekarang entri-entri pada kolom pertama di bawah posisi pivot bernilai nol. Selanjutnya, kita mencari kolom pivot kedua, yang dalam hal ini adalah kolom ketiga.
Di sini, $1$ pada baris kedua dan kolom ketiga berada pada posisi pivot.
Kita hanya perlu melakukan satu operasi baris untuk menghasilkan nol di bawah $1$.

Mengambil $-1$ kali baris kedua dan
menambahkannya ke baris ketiga menghasilkan
\begin{equation*}
\leftB
\begin{array}{rrr|r}
2 & 4 & -3 & -1 \\
0 & 0 &  1 & 1 \\
0 & 0 &  0 & 20
\end{array}
\rightB 
\end{equation*}

Kita dapat melanjutkan algoritme untuk membawa matriks ini ke \ef \;atau \rref. Namun, ingatlah bahwa
kita sedang mencari solusi sistem persamaan. Perhatikan kembali baris ketiga matriks tersebut.
Baris itu bersesuaian dengan
persamaan
\begin{equation*}
0x+0y+0z=20
\end{equation*}
Persamaan ini tidak memiliki solusi karena untuk semua $x,y,z$, ruas kiri sama dengan $0$ dan $0\neq 20.$ Hal ini menunjukkan bahwa
sistem persamaan yang diberikan tidak memiliki solusi. Dengan kata lain, sistem ini tak konsisten.
\end{solution}

Berikut adalah contoh lain untuk menemukan solusi sistem persamaan dengan membawa
matriks diperbesar yang bersesuaian ke \rref.

\begin{example}{Himpunan Solusi Tak Berhingga}\label{exa:infinitesetofsoln}
Berikan solusi lengkap sistem persamaan berikut \:
\begin{equation}
\begin{array}{ccccccccc}
3x&-&y&-&5z&=&9 \\
&&y&-&10z&=&0 \\
-2x&+&y&&&=&-6
\end{array}
\end{equation}
\end{example}

\begin{solution} Matriks diperbesar sistem ini adalah
\begin{equation*}
\leftB
\begin{array}{rrr|r}
 3 & -1 &  -5 & 9 \\
 0 &  1 & -10 & 0 \\
-2 &  1 &   0 & -6
\end{array}
\rightB
\end{equation*}
Untuk menemukan solusi sistem ini, kita akan membawa matriks diperbesar ke \rref{} dengan menggunakan Algoritme \ref{algo:rrefalgorithm}.
Kolom pertama merupakan kolom pivot pertama. Kita ingin menggunakan operasi baris untuk menghasilkan nol di bawah entri pertama dalam kolom ini, yang berada
pada posisi pivot pertama.
Ganti baris ketiga dengan $2$ kali baris pertama ditambah $3$ kali
baris ketiga. Hal ini menghasilkan
\begin{equation*}
\leftB
\begin{array}{rrr|r}
3 & -1 & -5 & 9 \\
0 & 1 & -10 & 0 \\
0 & 1 & -10 & 0
\end{array}
\rightB 
\end{equation*}
Sekarang kita telah menghasilkan nol di bawah $3$ dalam kolom pertama, sehingga kita beralih ke kolom pivot kedua (yaitu kolom kedua) dan mengulangi prosedurnya.
Ambil $-1$ kali baris kedua dan tambahkan ke baris ketiga.
\begin{equation*}
\leftB
\begin{array}{rrr|r}
3 & -1 &  -5 & 9 \\
0 &  1 & -10 & 0 \\
0 &  0 &   0 & 0
\end{array}
\rightB
\end{equation*}
Entri di bawah posisi pivot dalam kolom kedua kini bernilai nol. Perhatikan bahwa tidak ada lagi kolom pivot karena kita
hanya memiliki dua entri utama.

%%Solution using Style B definition of row-echelon form
Pada tahap ini, kita juga ingin agar entri utama sama dengan satu. Untuk melakukannya, bagi baris pertama dengan $3$.
\begin{equation*}
\leftB
\begin{array}{rrr|r}
1 & -\vspace{0.05in}\frac{1}{3} &  -\vspace{0.05in}\frac{5}{3} & 3 \\
0 &  1 & -10 & 0 \\
0 &  0 &   0 & 0
\end{array}
\rightB
\end{equation*}

Matriks ini kini berada dalam \ef.

Mari kita melanjutkan operasi baris hingga matriks berada dalam \rref.
Hal ini melibatkan pembuatan nol di atas posisi pivot pada setiap kolom pivot.
Proses ini hanya memerlukan satu langkah, yaitu menambahkan $ \frac{1}{3}$ kali baris kedua ke baris pertama.
\begin{equation*}
\leftB
\begin{array}{rrr|r}
1 & 0 & -5 & 3 \\
0 & 1 & -10 & 0 \\
0 & 0 & 0 & 0
\end{array}
\rightB 
\end{equation*}
%end of Kuttler specific solution 

%%Solution using original Kuttler definition of echelon form
%%This matrix is now in \ef. We want to continue with row operations
%%until the matrix is in \rref. Therefore, we can begin creating zeros above the
%%pivot positions in each pivot column. This requires only one step, which
%%is to take the second row and add it to the first row.
%%At this stage, we also want the leading entries to be equal to one. 
%%To do so, divide the first row by $3.$ These steps yield the matrix
%\begin{equation*}
%\leftB
%\begin{array}{rrr|r}
%%1 & 0 & -5 & 3 \\
%%0 & 1 & -10 & 0 \\
%%0 & 0 & 0 & 0
%\end{array}
%\rightB 
%\end{equation*}

Matriks ini berada dalam \rref; Anda sebaiknya memverifikasinya dengan menggunakan Definisi \ref{def:rref}.
 Persamaan-persamaan yang bersesuaian dengan \rref \;ini adalah
\[
\begin{array}{cccccccc}
x &-& &&5z&=&3 \\
&&y &-& 10z &=& 0 
\end{array}
\]
atau
\[
\begin{array}{ccccccc}
x&=&3&+&5z \\
y &=& &&10z
\end{array}
\]

Perhatikan bahwa $z$ tidak dibatasi oleh persamaan mana pun. Bahkan,
 $z$ dapat sama dengan sembarang bilangan. Sebagai contoh, kita dapat menetapkan $z = t$, dengan $t$ dapat dipilih sebagai sembarang bilangan.
Dalam konteks ini, $t$ disebut \textbf{parameter} \index{parameter}. Oleh karena itu, himpunan
solusi sistem ini adalah
\[
\begin{array}{c}
x=3+5t \\
y=10t \\
z=t
\end{array}
\]
dengan $t$ sembarang. Sistem tersebut memiliki himpunan solusi tak berhingga yang diberikan oleh persamaan-persamaan ini. Untuk setiap nilai $t$ yang kita pilih,
$x, y,$ dan $z$ diberikan oleh persamaan di atas. Sebagai contoh, jika kita memilih $t=4$, solusi yang bersesuaian adalah
\[
\begin{array}{c}
x = 3 + 5 (4) = 23\\
y = 10(4)=40 \\
z=4
\end{array}
\]
\end{solution}

Pada Contoh \ref{exa:infinitesetofsoln}, solusi melibatkan satu parameter. Solusi
suatu sistem dapat melibatkan lebih dari satu parameter, seperti diperlihatkan dalam contoh berikut.

\begin{example}{Himpunan Solusi dengan Dua Parameter}\label{exa:twoparametersetofsoln}
Temukan solusi sistem
\begin{equation*}
\begin{array}{ccccccccc}
x&+&2y&-&z&+&w&=&3 \\
x&+&y&-&z&+&w&=&1 \\
x&+&3y&-&z&+&w&=&5
\end{array}
\end{equation*}
\end{example}

\begin{solution} Matriks diperbesarnya adalah
\begin{equation*}
\leftB
\begin{array}{rrrr|r}
1 & 2 & -1 & 1 & 3 \\
1 & 1 & -1 & 1 & 1 \\
1 & 3 & -1 & 1 & 5
\end{array}
\rightB 
\end{equation*}
Kita ingin membawa matriks ini ke \ef. Di sini, kita akan menguraikan operasi baris yang digunakan. Namun, pastikan bahwa Anda
memahami langkah-langkah tersebut berdasarkan Algoritme \ref{algo:rrefalgorithm}.

Ambil $-1$ kali baris pertama dan tambahkan ke baris kedua. Kemudian ambil $-1$ kali
baris pertama dan tambahkan ke baris ketiga. Hal ini menghasilkan
\begin{equation*}
\leftB
\begin{array}{rrrr|r}
1 & 2 & -1 & 1 & 3 \\
0 & -1 & 0 & 0 & -2 \\
0 & 1 & 0 & 0 & 2
\end{array}
\rightB
\end{equation*}

%%Solution using Style B definition of row-echelon form
Sekarang tambahkan baris kedua ke baris ketiga dan bagi baris kedua dengan $-1$.
\begin{equation}
\leftB
\begin{array}{rrrr|r}
1 & 2 & -1 & 1 & 3 \\
0 & 1 & 0 & 0 & 2 \\
0 & 0 & 0 & 0 & 0
\end{array}
\rightB  \label{twoparameters1}
\end{equation}

%%Solution using original Kuttler definition of echelon form
%%Now add the second row to the third row.
%\begin{equation}
%\leftB
%\begin{array}{rrrr|r}
%%1 & 2 & -1 & 1 & 3 \\
%%0 & -1 & 0 & 0 & -2 \\
%%0 & 0 & 0 & 0 & 0
%\end{array}
%\rightB  \label{twoparameters1}
%\end{equation}

Matriks ini berada dalam \ef, dan kita dapat melihat bahwa $x$ dan $y$
bersesuaian dengan kolom pivot, sedangkan $z$ dan $w$ tidak. Oleh karena itu, kita akan menetapkan parameter bagi variabel $z$ dan $w$.
Tetapkan parameter $s$ untuk $z$ dan parameter
$t$ untuk $w.$ Baris pertama kemudian menghasilkan persamaan $x+2y-s+t=3$, sedangkan baris kedua menghasilkan
persamaan $y=2$. Karena $y=2$, persamaan pertama menjadi $x+4-s+t=3$,
yang menunjukkan bahwa solusi diberikan oleh
\begin{equation*}
\begin{array}{c}
x=-1+s-t \\
y=2 \\
z=s \\
w=t
\end{array}
\end{equation*}
Solusi ini lazim ditulis dalam
bentuk
\begin{equation}
\leftB
\begin{array}{c}
x \\
y \\
z \\
w
\end{array}
\rightB =\leftB
\begin{array}{c}
-1+s-t \\
2 \\
s \\
t
\end{array}
\rightB   \label{twoparameters2}
\end{equation}

\end{solution}

Contoh ini memperlihatkan suatu sistem persamaan dengan himpunan solusi tak berhingga yang bergantung pada dua parameter.
Dalam kasus himpunan solusi tak berhingga, penjelasannya dapat menjadi lebih jelas jika kita terlebih dahulu
menempatkan matriks diperbesar dalam \rref, bukan hanya
\ef, sebelum mencoba menuliskan deskripsi solusinya.

Dalam langkah-langkah di atas, ini berarti kita tidak berhenti pada \ef \;dalam persamaan \ref{twoparameters1}.
Sebaliknya, kita terlebih dahulu menempatkannya dalam \rref \;sebagai berikut.
\begin{equation*}
\leftB
\begin{array}{rrrr|r}
1 & 0 & -1 & 1 & -1 \\
0 & 1 & 0 & 0 & 2 \\
0 & 0 & 0 & 0 & 0
\end{array}
\rightB 
\end{equation*}
Selanjutnya, solusi dari baris kedua adalah $y=2$, sedangkan dari baris
pertama adalah $x=-1+z-w$. Jadi, dengan menetapkan $z=s$ dan $w=t,$ solusi diberikan oleh \ref{twoparameters2}.

Di sini Anda dapat melihat bahwa terdapat dua jalan menuju jawaban yang benar, dan keduanya menghasilkan
jawaban yang sama. Jadi, salah satu pendekatan dapat digunakan. Proses yang pertama kali kita gunakan dalam solusi di atas disebut \textbf{eliminasi Gauss}
\index{eliminasi Gauss}
Proses ini melibatkan membawa matriks ke \ef, mengubahnya kembali menjadi persamaan, dan menggunakan substitusi balik untuk menemukan solusi.
Ketika Anda melakukan operasi baris hingga memperoleh \rref, proses tersebut
disebut \textbf{eliminasi Gauss--Jordan}. \index{eliminasi Gauss--Jordan}

Kini kita telah menemukan solusi bagi sistem persamaan tanpa solusi dan dengan tak berhingga banyak solusi, baik dengan satu parameter maupun dua parameter.
Ingat kembali tiga jenis himpunan solusi yang dibahas pada bagian sebelumnya: tidak ada solusi, satu solusi, dan tak berhingga banyak solusi. Setiap
jenis solusi ini dapat dikenali dari grafik sistem. Ternyata kita juga dapat mengenali jenis solusi
dari \rref \;matriks diperbesar.

\begin{itemize}
\item {\em Tidak Ada Solusi: \em}
Jika sistem persamaan tidak memiliki solusi, \ef \;dari matriks diperbesar akan memiliki baris berbentuk
\[
\leftB
\begin{array}{rrrrr}
0 & 0 & 0 & | & 1
\end{array}
\rightB
\]

Baris ini menunjukkan bahwa sistem tersebut tak konsisten dan tidak memiliki solusi.

\item {\em Satu Solusi: \em}
Jika sistem persamaan memiliki satu solusi, setiap kolom
matriks koefisien merupakan kolom pivot.
Berikut adalah contoh matriks diperbesar dalam \rref \hspace{0.5mm} untuk sistem persamaan dengan satu solusi.
\begin{equation*}
\leftB
\begin{array}{rrr|r}
1 & 0 & 0 & 5 \\
0 & 1 & 0 & 0 \\ 
0 & 0 & 1 & 2 
\end{array}
\rightB
\end{equation*}

\item {\em Tak Berhingga Banyak Solusi: \em}
Jika sistem persamaan memiliki tak berhingga banyak solusi, solusinya memuat parameter. Akan
terdapat kolom-kolom matriks koefisien yang bukan kolom pivot.
Berikut adalah contoh matriks diperbesar dalam \rref \;untuk sistem persamaan dengan tak berhingga banyak solusi.
\begin{equation*}
\leftB
\begin{array}{rrr|r}
1 & 0 & 0 & 5 \\
0 & 1 & 2 & -3 \\ 
0 & 0 & 0 & 0 
\end{array}
\rightB
\end{equation*}
atau
\begin{equation*}
\leftB
\begin{array}{rrr|r}
1 & 0 & 0 & 5 \\
0 & 1 & 0 & -3
\end{array}
\rightB
\end{equation*}

\end{itemize}
